\section{Conclusion}
Implementing public policies demands various types of resources, such as common goods and services. It is crucial that the process of purchasing these goods and services be well planned and executed so that the government may use its budget efficiently and achieve public policy goals effectively.

This paper investigates the enforcement costs of health litigation and administrative requests for the public budget. We evaluate the government waste generated when the judiciary directly affects public policy. In this case, health litigation imposes multifold restrictions on the public procurement process, harming tender outcomes.

From a policy perspective, this research indicates that judges should consider the public budget implications and administrative costs of purchasing health items under pressure in their decisions. An institutional arrangement integrating the judiciary and the executive branches, enabling joint actions, might mitigate waste in the health litigation context.
